\section{Geometric recovery training and correction}\label{app:geometric-recovery}
\paragraph{Backbone continuation.}
Each initial FM fit receives 4,000 additional AdamW updates at learning rate $10^{-5}$ and batch size 32, with gradient-norm limit one and EMA decay capped at 0.999. Odd updates use the original two-pass FM loss. Even updates draw progress uniformly from $[0.60,0.95]$ and regress both endpoint predictions to the paired OMol25 reference. The continuation control shares these batches, progress values, and update counts but omits added perturbations and local geometry losses.

Recovery updates perturb the late interpolant. Each atom receives Gaussian noise independently with probability 0.25; its standard deviation is chosen uniformly from 0.05, 0.15, and 0.30\,\AA\ per molecule. With probability 0.5, a random spatial half of the molecule also receives a shared Gaussian displacement with standard deviation 0.35\,\AA. We subtract the mean displacement and retain the original reference as the target.

The recovery loss weights endpoint mean squared error by four, local distance error by 0.5, and local angle error by 0.1. For reference distance $d_{ij}$ and radius sum $\ell_{ij}$, distinct-atom pairs have weight $w_{ij}=\operatorname{sigmoid}[(1.30-d_{ij}/\ell_{ij})/0.10]$. Distance errors use distances normalized by $\ell_{ij}$; angle errors compare neighbor-direction cosines with weight $w_{ij}w_{ik}$ for distinct $i,j,k$. Each local loss is normalized by its weight sum within a molecule, then averaged over molecules and both passes. All targets are derived from reference coordinates rather than bond labels.

\paragraph{Geometric-head targets.}
A fixed hash order selects 4,096 of the 20,000 training rows; 3,783 pass the graph, connectivity, overlap, zero-assigned-radical, and geometry screen. Selection outcomes are recorded, and evaluation compositions are excluded. Inferred graphs are used for screening only; network inputs are coordinates, elements, and progress.

We draw $p$ uniformly from $[0.4,0.95]$ and perturb the reference with centered Gaussian noise of scale $s(p)=0.06+0.30(1-p)$\,\AA\ per coordinate. Fifteen percent of examples retain the original coordinates and receive zero targets. The remaining targets are $(X_{\rm ref}-\widetilde H)/(1-p)$, uniformly rescaled to cap the largest atomic norm at $4p^2$\,\AA\ per unit progress while preserving zero total velocity. Squared vector errors have weight 0.25 for hydrogen and one for heavy atoms, normalized by the molecular weight sum. Training uses 10,000 AdamW updates at learning rate $10^{-3}$ and batch size 32, gradient-norm limit one, and EMA decay capped at 0.995. Sampling uses the final EMA.

\paragraph{Geometric-head architecture.}
The 7,747-parameter head contains a 16-dimensional element embedding, a learned linear neighbor-embedding update, and two width-64 SiLU layers. Contact weights are $\operatorname{sigmoid}[(1.30-d_{ij}/\ell_{ij})/0.20]$. Each node combines its element embedding with contact-weighted neighbor embeddings. Symmetric pair sums and products are concatenated with distances, coordination summaries, six radial basis functions, the radius sum, and $p,p^2$.

Three bounded coefficients scale the normalized relative vector and two bounded radial variants. Each pair contributes opposite vectors to its endpoints. Normalizing by three times the maximum weighted degree, floored at one, and multiplying by $4p^2$ bounds each atomic correction by $4p^2$. The final layer starts at zero. Tests check rotation, reflection, permutation, and centering. A parameter-matched variant with directional moments and nonradial vectors did not outperform the radial head during development.

\paragraph{Combined correction and cost.}
The geometric head first updates the endpoint at strength one. The physical head then uses that endpoint and the current trajectory state at strength four. Sampling fixes both heads and the backbone, using 128 backbone calls and 64 calls to each head. We reuse the original 7,106-parameter physical head, so geometric refinement requires no additional eSEN labels. Per fit, recovery adds 256,000 backbone example-forwards, while geometric-head training adds 320,000 small-network example-forwards and no backbone evaluations.

\paragraph{Development.}
We compare recovery and geometric-head variants on 24 training-disjoint validation compositions, split equally between 17--28 and 29--40 atoms. Each of two fits generates 16 samples per composition. Table~\ref{tab:geometry-development} reports the development sequence. We select the radial head and recovery with local geometry supervision, then fix all choices before generating new outputs on the existing 64-composition benchmark. That benchmark had already been used in earlier project work.

\begin{table}[htbp]\centering\small
\caption{Geometric-recovery development results, with 768 raw outputs per row. The last column additionally requires GFN2 force RMS at most 5 eV/\AA.}
\begin{tabular}{p{.49\linewidth}rrr}\toprule
Setting & Graph & Geometry & Geom. + force\\\midrule
Initial FM + physical head &32.55&9.11&8.85\\
Continuation control + physical head &36.07&10.94&10.55\\
Perturbed recovery + physical head &41.67&10.94&10.68\\
Recovery with local geometry + physical head &43.62&12.11&11.85\\
Initial FM + radial geometry + physical head &36.33&11.33&11.33\\
Initial FM + moment geometry + physical head &36.98&10.94&10.94\\
Recovery with local geometry + both heads &44.40&14.45&14.32\\\bottomrule
\end{tabular}\label{tab:geometry-development}\end{table}

On the development panel, the combined method adds 5.47 geometry-and-force points over the original model, with fit-specific gains of 4.43 and 6.51. The descriptive composition interval is $[1.82,9.38]$ and the crossed fit/composition interval is $[1.04,10.03]$. Incremental comparisons against the individual components have wider intervals containing zero.

\begin{table}[htbp]\centering\small
\caption{Geometric and physical corrections on shared recovery-trained parents. Percentages count all 768 attempts per row.}
\begin{tabular}{lrrr}\toprule
Correction & Geometry & Geom. + force $\leq2$ & Geom. + force $\leq5$\\\midrule
None &10.68&0.13&6.77\\
Physical &12.11&1.56&11.85\\
Geometric &14.19&1.82&13.54\\
Geometric + physical &14.45&3.13&14.32\\\bottomrule
\end{tabular}\label{tab:geometric-physical-fields}\end{table}

On the 64-composition benchmark, the first two fits improve geometry-and-force yield from 4.88\% to 10.25\%, compared with 5.76\% for matched-update continuation. Each row contains 2,048 attempts. Gains over the original and continued models are 5.37 points ($[3.56,7.13]$) and 4.49 points ($[2.44,6.45]$), respectively, with both fits improving. The controls retain their physical head; the complete method additionally uses the geometric head. These comparisons therefore evaluate the full refinement procedure. In a separate 128-output fixed-noise diagnostic, increasing backbone calls from 128 to 512 changed graph-valid successes from 38 to 39 and left geometry-and-force successes at 13.

\paragraph{Additional fits.}
We apply the selected recovery and geometric-head procedure to three additional FM parents, retaining their physical heads. Geometry-and-force successes rise from 22, 54, and 51 to 61, 101, and 121 out of 1,024 per fit. The gains are 3.81, 4.59, and 6.84 points, with mean 5.08 and crossed 95\% interval $[3.09,7.52]$. Across five fits, 493 of 5,120 outputs pass geometry and force, compared with 227 for the initial model. Geometry alone accepts 494 outputs, all with distinct inferred connectivities within each composition and fit.
